\documentclass[10pt]{article}

\usepackage[margin=1in]{geometry}
\usepackage{amsmath,amssymb,booktabs,graphicx,microtype}
\usepackage[numbers,sort&compress]{natbib}
\usepackage[hidelinks]{hyperref}
\usepackage{xcolor}

\newcommand{\placeholder}[1]{\textcolor{gray}{\textit{#1}}}
\newcommand{\paperidea}[1]{\item #1}

\title{How Many Bits Does Learning Require?\\
\large Conditional Description Length of Learned LLM Behavior}
\author{Joey David \and \placeholder{coauthors}}
\date{}

\begin{document}
\maketitle

\begin{abstract}
\begin{itemize}
  \paperidea{One-sentence problem: after a pretrained model acquires a capability through fine-tuning, how many transmitted bits are actually required to preserve that behavioral change?}
  \paperidea{Treat the frozen base model and adapter architecture as shared side information; measure the exact decoder-sufficient description length of the learned update.}
  \paperidea{Main empirical headline to fill after replication: very low-rate descriptions retain a large fraction of downstream task gain, while nominal precision alone is a poor predictor of behavior.}
  \paperidea{Main methodological contribution: exact serialized-rate accounting plus adaptive variable-length coding (conditional-MDL view), evaluated after an actual encode/decode round trip.}
  \paperidea{Reasoning extension: compare the information budget needed to preserve final-answer capability with the budget needed to preserve the model's reasoning trajectory.}
\end{itemize}
\end{abstract}

\section{Introduction}
\begin{itemize}
  \paperidea{Motivation: fine-tuning writes new behavior into a pretrained system; adapter parameter count is only a loose upper bound on how much information was actually written.}
  \paperidea{Communication view: training data $\rightarrow$ learned update $\Delta\theta$ $\rightarrow$ transmitted description $z$ $\rightarrow$ reconstructed update $\widehat{\Delta\theta}$ $\rightarrow$ behavior.}
  \paperidea{Distinguish three quantities clearly: nominal bit width, exact artifact description length, and behavioral information / retained task gain.}
  \paperidea{Key question 1: what is the rate--performance curve for learned behavior?}
  \paperidea{Key question 2: does allocating bits non-uniformly across an adapter move this frontier?}
  \paperidea{Key question 3 (reasoning): is preserving a correct answer cheaper than preserving the computation / trajectory that produced it?}
  \paperidea{Potential contributions list: exact coding protocol; cross-task rate--performance curves; evidence that weight reconstruction error does not fully predict behavior; adaptive MDL frontier; outcome-vs-trajectory reasoning analysis.}
  \paperidea{Do not sell this as ``LoRA compression'' alone; emphasize learned-information measurement and rate--distortion of acquired behavior.}
\end{itemize}

\section{Related Work}
\subsection{Parameter-efficient adaptation and learned deltas}
\begin{itemize}
  \paperidea{LoRA as the learned channel: cite \citet{hu2021lora}.}
  \paperidea{QLoRA for low-precision base-model adaptation context: \citet{dettmers2023qlora}.}
  \paperidea{Delta tuning / parameter-efficient adaptation as the broader family.}
\end{itemize}

\subsection{Compression of fine-tuning updates}
\begin{itemize}
  \paperidea{BitDelta: one-bit full-model fine-tuning deltas; closest conceptual precedent for ``fine-tuning writes surprisingly little information'' \citep{liu2024bitdelta}.}
  \paperidea{Delta-CoMe: mixed-precision compression of task-specific deltas \citep{ping2024deltacome}.}
  \paperidea{LoRAQuant: directly overlapping mixed-precision LoRA compression; distinguish our objective (exact description-length / behavioral frontier and adaptive coding) from a deployment-compression objective \citep{mirzaei2025loraquant}.}
  \paperidea{LoRA-Squeeze: post-hoc rank compression; useful comparison between restricting the channel during training and compressing an already learned solution \citep{vulic2026lorasqueeze}.}
\end{itemize}

\subsection{Minimum description length and information in weights}
\begin{itemize}
  \paperidea{Historical anchor: minimizing the description length of network weights \citep{hinton1993mdl}.}
  \paperidea{Be precise: our main measured quantity is a concrete conditional code length, not a claim that the current codec is the globally shortest possible description.}
  \paperidea{Suggested framing: empirical upper bounds on $L(\Delta\theta\mid M,\mathcal A,\mathcal D)$ under a specified decoder family.}
\end{itemize}

\subsection{Reasoning processes and monitorability}
\begin{itemize}
  \paperidea{Process supervision / step-level reasoning as a distinct target from outcome accuracy \citep{lightman2023verify}.}
  \paperidea{Chain-of-thought may be behaviorally useful without being fully faithful; motivate measuring latent/conditional trajectory preservation rather than treating surface text as ground truth \citep{lyu2023faithful}.}
  \paperidea{Safety relevance of preserving / monitoring reasoning traces \citep{korbak2025monitorability}.}
\end{itemize}

\section{Problem Formulation}
\begin{itemize}
  \paperidea{Frozen pretrained model $M$; learned LoRA tensors $\Delta\theta$; codec/decoder family $\mathcal D$.}
  \paperidea{Conditional description length: $R = |z|$ bits for a self-contained code $z$ that reconstructs $\widehat{\Delta\theta}$ given shared $M$, adapter architecture, and decoder.}
  \paperidea{Report both total bits and effective bits per original learned scalar: $R/N$.}
  \paperidea{Task distortion should be defined behaviorally, e.g. retained gain $(S(\widehat{\Delta\theta})-S(M))/(S(\Delta\theta)-S(M))$, with uncapped values retained in tables and capped values only for frontier visualization.}
  \paperidea{Separate weight-space distortion (relative RMSE) from behavioral distortion. One central analysis should test how tightly they correlate.}
  \paperidea{Rate--distortion / Pareto definition: lower exact description length and higher retained behavioral gain.}
\end{itemize}

\section{Experimental Protocol}
\begin{itemize}
  \paperidea{Frozen bases: Mistral-7B and Qwen2.5-7B; pin revisions.}
  \paperidea{Main adapter: rank-16 full LoRA over attention and MLP projections; preserve matched raw checkpoint across every codec.}
  \paperidea{Tasks: Magicoder $\rightarrow$ HumanEval; MetaMathQA $\rightarrow$ GSM8K/MATH; XSum $\rightarrow$ XSum.}
  \paperidea{Three seeds per main cell; report all cells, including failures / no-learning gates.}
  \paperidea{Raw learning gate: show held-out NLL improvement before making any compression claim.}
  \paperidea{Exact accounting: headers, selectors, scales, packed values, padding, entropy compression, final file size.}
  \paperidea{Every evaluated point must be decoded from the actual stored artifact before task evaluation.}
\end{itemize}

\section{Fixed-Precision Rate--Performance Curves}
\subsection{Code families}
\begin{itemize}
  \paperidea{Define in plain language, not codec jargon: sign-only representation; four-level sign-preserving representation; progressively finer symmetric representations.}
  \paperidea{Main family (mid-rise): zero-free symmetric levels $\pm1,\pm3,\dots$ with a least-squares row scale, so every codeword is used and no clipping percentile has to be tuned. At one bit this is exactly $\mathrm{sign}(w)\cdot\mathrm{mean}(|w|)$ per row.}
  \paperidea{Control family (mid-tread): absmax scale with an exact zero level, the default geometry in most quantization code. At two bits it represents only $-1,0,+1$, spending a quarter of its codebook on an unused symbol.}
  \paperidea{Report both at matched \emph{payload} width and at exact file rate; they differ, because a three-symbol alphabet leaves more for the entropy coder to remove. This is itself an argument for exact file accounting.}
  \paperidea{Include FP16 / 8 / 4 / 3 / 2 / 1-bit controls and LoRAQuant controls.}
\end{itemize}

\subsection{Preliminary observations to verify across the full campaign}
\begin{itemize}
  \paperidea{Magicoder/Mistral seed-33: raw LoRA improves HumanEval from 11/164 to 26/164.}
  \paperidea{Sign-only adapter retains 23/164 at roughly 1.46 effective bits per learned scalar.}
  \paperidea{A sign-preserving four-level representation reaches approximately raw performance at roughly 2.36 effective bits per learned scalar.}
  \paperidea{Weight-space calibration of the geometry claim, measured directly and independent of any task: at a matched payload width the mid-tread code costs about $2.5\times$ the relative reconstruction error of the mid-rise code at two bits, $1.55\times$ at three, and $1.07\times$ at four. The gap is a low-bit phenomenon and closes as width grows.}
  \paperidea{The two-bit mid-tread code lands near the one-bit mid-rise file rate yet reconstructs worse, so it is not merely dominated, it is dominated by a code with half its nominal width. This is the cleanest available evidence that ``number of bits'' alone does not characterize a code.}
  \paperidea{Caution: the earlier campaign used the mid-tread geometry as its \emph{default} uniform codec, which made its 2- and 3-bit points, and hence its reported 2-to-3-bit transition, an artifact of the codec rather than a property of the adapter. Re-derive every fixed-rate point before reusing any of those numbers.}
  \paperidea{Non-monotonic HumanEval results across 2/3/4/8-bit points: do not overinterpret until replicated; compare with held-out NLL and uncertainty across seeds.}
\end{itemize}

\subsection{Figures and tables}
\begin{itemize}
  \paperidea{Figure 1: exact effective bits/scalar vs retained task gain, all codecs, with nondominated frontier.}
  \paperidea{Figure 2: relative weight RMSE vs retained task gain to show when numerical reconstruction diverges from behavioral preservation.}
  \paperidea{Table 1: raw/base/task scores, exact file bytes, effective rate, retained gain, RMSE, seed.}
\end{itemize}

\section{Adaptive Description Length / MDL Experiment}
\begin{itemize}
  \paperidea{Each LoRA row chooses among: omit it entirely, or encode it with 1/2/3/4/8-bit sign-preserving quantization. The row quantizer is the same one used for the fixed-rate curves, so the adaptive and fixed families differ only in allocation, not in code geometry.}
  \paperidea{Current allocation objective: per-row reconstruction error $+\lambda\times$ proxy code length. Sweep $\lambda$ / target budgets; plot actual final artifact bits, not proxy rate.}
  \paperidea{Make the conditional-MDL qualification explicit: base model, architecture, and decoder are shared side information; the code is a concrete upper bound, not proof of globally minimal Kolmogorov/MDL description.}
  \paperidea{Main comparison: adaptive allocation vs best fixed-precision codec at matched exact file rate.}
  \paperidea{Report allocation histograms: fraction of rows and fraction of values receiving 0/1/2/3/4/8 bits.}
  \paperidea{Dense follow-up sweep only around the knee of the first coarse frontier.}
  \paperidea{Stronger extension: replace weight-MSE importance with calibration-NLL sensitivity, Fisher/Hessian approximations, or direct behavioral loss; ask whether behavior-aware allocation shifts the frontier.}
  \paperidea{Ablation: row-wise vs tensor-wise vs rank-component-wise allocation, including metadata cost.}
\end{itemize}

\section{Where Are the Learned Bits?}
\begin{itemize}
  \paperidea{Layer/projection decomposition of selected bit budgets.}
  \paperidea{Attention vs MLP matched-parameter placement controls already specified in the experiment plan.}
  \paperidea{LoRA-A vs LoRA-B asymmetry; rank-component concentration; whether high-bit regions recur across seeds/tasks.}
  \paperidea{Cross-seed overlap: do independently trained adapters spend their scarce bits in the same modules?}
  \paperidea{Potential result language: ``behaviorally relevant information is non-uniformly distributed'' only if adaptive coding consistently dominates matched-rate uniform coding.}
\end{itemize}

\section{Reasoning: Outcome Bits vs Trajectory Bits}
\subsection{Core experiment}
\begin{itemize}
  \paperidea{Train a reasoning adapter on explicit worked solutions (MetaMathQA or a cleaner curated subset); evaluate on GSM8K/MATH.}
  \paperidea{Generate/fix a raw-adapter reasoning trajectory for each evaluation problem.}
  \paperidea{For every compressed adapter, teacher-force that exact trajectory and measure token-level divergence from the raw adapter along reasoning tokens.}
  \paperidea{Primary continuous trajectory metric: average $D_{\mathrm{KL}}(p_{\mathrm{raw}}(\cdot\mid x,r_{<t})\Vert p_{\mathrm{compressed}}(\cdot\mid x,r_{<t}))$ on reasoning positions.}
  \paperidea{Primary behavioral metric: final-answer correctness. Plot both against exact adapter description length.}
  \paperidea{Central question: does trajectory fidelity collapse at a higher rate than final-answer accuracy?}
\end{itemize}

\subsection{Answer-only vs chain-of-thought training}
\begin{itemize}
  \paperidea{Matched problems, two adapters: problem$\rightarrow$answer only vs problem$\rightarrow$reasoning$\rightarrow$answer.}
  \paperidea{Estimate $R^*_{\mathrm{answer}}(a)$ and $R^*_{\mathrm{CoT}}(a)$: minimum description length needed to retain target accuracy $a$.}
  \paperidea{Interpret their difference as an operational ``information cost of explicit reasoning'', with caveats about optimization / training-format differences.}
\end{itemize}

\subsection{Latent trajectory extension}
\begin{itemize}
  \paperidea{At sentence/step boundaries, compare hidden states of raw and compressed adapters (cosine / CKA / normalized distance).}
  \paperidea{Reuse reasoning-trajectory work: solution-object updates, step transitions, and layer-local state changes.}
  \paperidea{Ask which internal transitions disappear first as description length shrinks, and whether final accuracy survives after trajectory structure changes.}
\end{itemize}

\section{Ablations and Robustness}
\begin{itemize}
  \paperidea{Models, tasks, three seeds, adapter ranks, LoRA placement.}
  \paperidea{Alternative row scales / codebooks / entropy coding; show exact-file-rate matching.}
  \paperidea{No entropy compressor vs zlib to distinguish structural compressibility from symbol-frequency compressibility.}
  \paperidea{Calibration-set size for adaptive allocation.}
  \paperidea{Task-score confidence intervals / bootstrap over examples; do not treat 25 vs 27 HumanEval solves as a meaningful improvement without replication.}
  \paperidea{Held-out NLL as a higher-resolution companion metric to discrete benchmark accuracy.}
\end{itemize}

\section{Discussion}
\begin{itemize}
  \paperidea{What does a short adapter description mean about the amount of task-specific information written during fine-tuning?}
  \paperidea{Distinguish compressibility of the learned solution from simplicity of the training data / function.}
  \paperidea{Why exact code length is an upper bound, not an intrinsic invariant of the capability.}
  \paperidea{Implications for modular model storage are secondary; scientific focus is where and how much new information is needed to alter behavior.}
  \paperidea{Reasoning interpretation if outcome survives at lower rates than trajectory: many internal computations may be redundant with respect to the final task metric.}
\end{itemize}

\section{Limitations}
\begin{itemize}
  \paperidea{Codec-family dependence of measured description length.}
  \paperidea{LoRA-only conclusions do not immediately generalize to full fine-tuning.}
  \paperidea{Benchmark accuracy is a coarse behavioral distortion measure.}
  \paperidea{Teacher-forced trajectory KL measures preservation relative to one raw trajectory, not ground-truth reasoning faithfulness.}
  \paperidea{Finite models/tasks/seeds; avoid universal claims about ``bits required for learning''.}
\end{itemize}

\section{Conclusion}
\begin{itemize}
  \paperidea{Restate the paper as measurement of learned behavioral information, not merely adapter compression.}
  \paperidea{One-line final claim slot after replication.}
\end{itemize}

\bibliographystyle{plainnat}
\bibliography{references}

\end{document}
